\documentclass[11pt]{article}
\usepackage[T1]{fontenc}
\usepackage{lmodern,microtype,amsmath,amssymb,amsthm,booktabs,longtable,array}
\usepackage[margin=1in]{geometry}
\usepackage{xurl}
\usepackage[hidelinks]{hyperref}
\setlength{\emergencystretch}{3em}
\newtheorem{theorem}{Theorem}
\newtheorem{lemma}[theorem]{Lemma}
\newtheorem{proposition}[theorem]{Proposition}
\providecommand{\tightlist}{\setlength{\itemsep}{0pt}\setlength{\parskip}{0pt}}
\providecommand{\passthrough}[1]{#1}
\hypersetup{pdfauthor={Yongxian Zhang}}

\title{An orbital primal-dual theorem for finite permutation actions}
\author{Yongxian Zhang\\[2pt]\small School of Computer Science and Engineering\\\small South China University of Technology\\\small\texttt{mxymmxym1@gmail.com}\\\small ORCID: 0009-0000-3864-3536}
\date{8 October 2026}
\begin{document}\maketitle
\begin{abstract}For any finite group acting on a nonempty finite set, we identify the sharp atom response under marginal-preserving probability perturbations with a rational orbital linear program and its minimum-oscillation dual. Rational optimizers exist, and positive optima are attained by disjoint central laws and sharp small perturbations of the uniform law. Neither faithfulness nor transitivity is assumed. All principal conclusions have complete Lean proofs.\end{abstract}
\section{The exact finite-action
theorem}\label{the-exact-finite-action-theorem}

Let \(G\) be a finite group acting (not necessarily transitively or
faithfully) on a finite nonempty set \(\Omega\), with \(e\in G\) its
identity. Write \(u_G\) for the uniform law on \(G\), and let
\(\mathcal V\) be the real vector space of signed functions
\(v:G\to\mathbb R\) satisfying

\[
\sum_{g\in G}v(g)=0,\qquad
\sum_{\substack{g\in G\\g(x)=y}}v(g)=0
\quad\text{for all }x,y\in\Omega. \tag{1}
\]

Let \(\mathcal O_1,\ldots,\mathcal O_r\) be the orbitals of the action,
i.e.~the orbits of \(G\) on \(\Omega\times\Omega\) under simultaneous
relabeling. Define their \emph{orbital displacement counts}

\[
F_j(g)=\#\{x\in\Omega:(x,g(x))\in\mathcal O_j\}. \tag{2}
\]

Each \(F_j\) is constant on conjugacy classes of \(G\). Let
\(C_1,\ldots,C_s\) be those conjugacy classes, with \(C_1=\{e\}\), and
put

\[
M_{ji}=\frac1{|C_i|}\sum_{g\in C_i}F_j(g),\qquad
M\in\mathbb Q^{r\times s}. \tag{3}
\]

For \(\mathcal V\ne\{0\}\), define

\[
\mathcal C(G,\Omega)=\sup_{\substack{v\in\mathcal V\\v\ne0}}
\frac{|v(e)|}{\tfrac12\sum_g|v(g)|}. \tag{4}
\]

For \(\mathcal V=\{0\}\), set \(\mathcal C(G,\Omega)=0\). The supremum
is finite and \(0\le\mathcal C(G,\Omega)\le1\).

\textbf{Theorem (exact orbital primal-dual characterization).} The
following four quantities are all equal:

\[
\begin{aligned}
\mathcal C(G,\Omega)
&=\max_{\substack{P,Q\text{ probability laws on }G\\
P\{g:gx=y\}=Q\{g:gx=y\}\ \forall x,y}}
\big(P(e)-Q(e)\big)\\
&=\max_{\substack{p,q\in\mathbb R_{\ge0}^{s}\\
\mathbf1^Tp=\mathbf1^Tq=1,\ Mp=Mq}}
(p_1-q_1)\\
&=\min_{\lambda\in\mathbb R^r}
\left[
\max_{1\le i\le s}\big(\mathbf1_{\{i=1\}}-\lambda^TM_{\cdot i}\big)
-\min_{1\le i\le s}\big(\mathbf1_{\{i=1\}}-\lambda^TM_{\cdot i}\big)
\right]\\
&=\min_{\varphi\in\mathrm{span}_{\mathbb R}\{F_1,\ldots,F_r\}}
\operatorname{osc}_{g\in G}\big(\mathbf1_{\{g=e\}}-\varphi(g)\big).
\end{aligned}
\tag{5}
\]

Both finite LPs in the middle have \textbf{rational optimal solutions},
so \(\mathcal C(G,\Omega)\in\mathbb Q\). If \(\mathcal C(G,\Omega)>0\),
there exist two \textbf{conjugation-invariant, disjointly supported}
probability laws \(P,Q\) with identical one-point marginals and
\(P(e)-Q(e)=\mathcal C(G,\Omega)\). Hence for all sufficiently small
\(\delta\ge0\),

\[
\nu_\delta=u_G+\delta(P-Q) \tag{6}
\]

is a probability law with exactly the same one-point marginals as
\(u_G\) and

\[
\|\nu_\delta-u_G\|_{\mathrm{TV}}=\delta,\qquad
\nu_\delta(e)-u_G(e)=\mathcal C(G,\Omega)\delta. \tag{7}
\]

For an arbitrary probability law \(\nu\) whose marginals match \(u_G\),
and any \(\sigma\in G\),

\[
\left|\nu(\sigma)-\frac1{|G|}\right|
\le\mathcal C(G,\Omega)\|\nu-u_G\|_{\mathrm{TV}}, \tag{8}
\]

with this constant optimal whenever it is positive. Left translation
moves the sharp construction from \(e\) to any prescribed \(\sigma\).

\textbf{Proof.} First formulate a finite \textbf{unreduced} LP. Choose
\(P,Q\) as nonnegative probability vectors indexed by \(G\), impose
their equality of all one-point marginals, and maximize \(P(e)-Q(e)\).
Call its optimum \(\gamma\).

For any feasible pair, \(v=P-Q\in\mathcal V\),
\(\|v\|_{\mathrm{TV}}\le1\), so \(P(e)-Q(e)\le\mathcal C(G,\Omega)\).
Conversely, for any nonzero \(v\in\mathcal V\), its Jordan parts
\(v_+,v_-\) each have mass \(d=\|v\|_{\mathrm{TV}}>0\), and the
probability measures \(P=v_+/d,\ Q=v_-/d\) have identical marginals.
Exchanging them if necessary, they give objective \(|v(e)|/d\). Hence
\(\gamma=\mathcal C(G,\Omega)\). If \(\mathcal V=\{0\}\), the two
quantities are both zero.

The unreduced primal is a rational finite LP. Its dual has one
unrestricted multiplier \(\lambda_{xy}\) for each marginal equality and
two normalization multipliers \(a,b\). Write

\[
\phi(g)=\sum_{x,y\in\Omega}\lambda_{xy}\,\mathbf1_{\{g(x)=y\}}.
\]

Dual feasibility is exactly

\[
a\ge\mathbf1_{\{g=e\}}-\phi(g),\qquad
b\ge-\mathbf1_{\{g=e\}}+\phi(g)
\quad\text{for all }g\in G.
\]

Minimizing \(a+b\) for fixed \(\phi\) gives
\(\max_g(\mathbf1_{\{g=e\}}-\phi(g))-
\min_g(\mathbf1_{\{g=e\}}-\phi(g))\), its oscillation. The primal is
feasible (take \(P=Q=u_G\)) and bounded, so \textbf{finite-dimensional
LP strong duality} gives equality with the minimum oscillation. The LP
has rational coefficients and a finite attained optimum; standard
rational Gaussian elimination at a basic feasible optimum gives rational
primal and dual certificates.

Next average an optimal primal pair under simultaneous conjugation by
\(h\in G\): replace \(P,Q\) by the averages of their pushforwards under
\(g\mapsto hgh^{-1}\). Equal one-point marginals remain equal because
conjugating both pairs of image coordinates by \(h\) permutes the
constraints. The identity atom and the objective are unchanged. Thus a
conjugation-invariant optimal pair exists. Such pairs assign a constant
probability density to each conjugacy class, and their equality of
one-point marginals is equivalent to the reduced moment equation
\(Mp=Mq\): for a central law, each entry \(\Pr(gx=y)\) is constant on
the orbital containing \((x,y)\), and its orbital average is
\(\mathbb EF_j/|\mathcal O_j|\). This proves the second line of (5).

Similarly, average any dual function \(\phi\) over conjugations. The
indicator \(\mathbf1_{\{g=e\}}\) is conjugacy invariant, while
oscillation is convex and invariant under conjugation; the average
cannot increase oscillation. A conjugation-averaged linear combination
of the image indicators has coefficients constant on their simultaneous
\(G\)-orbits in \(\Omega\times\Omega\). Therefore it lies in
\(\mathrm{span}\{F_j\}\). This proves the fourth line of (5). Evaluating
the same oscillation on conjugacy classes gives the third line, whose
entries are precisely (3).

When the optimum \(\gamma>0\), an optimal reduced pair \(P,Q\) must be
\textbf{mutually singular}. Otherwise \(d=\|P-Q\|_{\mathrm{TV}}<1\), and
its Jordan parts normalized by \(d\) would form an admissible pair with
strictly larger objective \(\gamma/d>\gamma\), contradiction. Thus the
supports of an optimal central \(P,Q\) are disjoint and (6) is a
probability measure for every \(0\le\delta\le(|G|\max_gQ(g))^{-1}\),
with all the stated identities.

Finally, for any law \(\nu\) with the uniform law's one-point marginals,
\(v=\nu-u_G\in\mathcal V\) and the defining norm bound gives (8) at the
identity. Replacing \(\nu\) by its left translate by \(\sigma^{-1}\)
preserves marginal equality with \(u_G\), total variation and atom
excess, proving (8) for every \(\sigma\). The same translation
transports the attained equality. QED.

\section{Context and scope}\label{context-and-scope}

This theorem extracts the universal finite-action principle from Section
18 of the repository's sharp robust permanent note. The separate closed
formulas for particular symmetric-group actions are outside this paper.
The duality and rational-attainment inputs are classical
finite-dimensional linear programming; see A. Schrijver, \emph{Theory of
Linear and Integer Programming}, Wiley, 1986. The formal proof supplies
the required real-dual and rationalization bridges using proved
finite-dimensional linear algebra and Hahn--Banach. No claim of
resolving a famous open conjecture is made.

\section{Formal verification and reproduction}
The principal theorem has a complete Lean proof. The precise scope is:
Only the universal finite-action theorem, rational primal/dual attainment, all-atom sharpness and degenerate zero-optimum case. Separate explicit symmetric-group formulas in the parent note are not certified by this paper.
The public proof source snapshot is fixed at commit
\nolinkurl{02af7db99db789fecb0d5f5f8f37b1a4e40762a0} in \url{https://github.com/mxym/math}.
The proof packages, with their own dependency locks, audit records and
reproduction instructions, are:
\begin{itemize}
\item \nolinkurl{formalizations/orbital-primal-dual}.
\end{itemize}
The main formal entry points are:
\begin{itemize}
\item \nolinkurl{OrbitalMarginals.universal_orbital_theorem}.
\item \nolinkurl{OrbitalMarginals.orbital_span_exact_minimum}.
\item \nolinkurl{OrbitalMarginals.sharp_attainment_all_atoms}.
\end{itemize}
Lean is fixed to version 4.34.1, compiler commit
\nolinkurl{5045d0056413266e57c625dcd7c365b10e377c52}; Mathlib is fixed to
\nolinkurl{d13f23b723b8a846827a245b89c10fc7d3f11612}.
The audited mathematical roots depend only on the standard axioms
\texttt{propext}, \texttt{Classical.choice}, and \texttt{Quot.sound}.
The archived audits state their actual compilation and replay scopes;
publication of this typeset version does not imply a new execution of every
previous empty-kernel replay. The companion README maps each printed
principal result to its formal theorem and explains how to repeat verification.

\section*{Research and preparation disclosures}
This research, formalization and manuscript preparation used AI assistance.
The author is responsible for the manuscript. The principal claims are
supported by the readable proofs and the public Lean artifacts specified
above. Formal verification is not external human peer review; no such review
is claimed. No external funding was received. No worldwide priority or
journal acceptance is asserted.
\end{document}
